\documentclass[10pt,a4paper]{amsart}
\usepackage[margin=19mm]{geometry}
\usepackage{amsmath,amssymb,mathtools}
\usepackage[scr=rsfs,cal=euler]{mathalfa}
\usepackage{xfrac}
\usepackage[unicode,hidelinks,bookmarks=false]{hyperref}
\newcommand{\ie}{i.e.\ }
\newcommand{\cf}{cf.\ }
\newcommand{\resp}{respectively\ }
\newcommand{\Subsection}{\subsection}
\providecommand{\nobreakdash}{\nobreak\hbox{-}}
\setlength{\emergencystretch}{2em}
\multlinegap0pt
\usepackage{bm}
\usepackage{dsfont}
\usepackage{amssymb}
\usepackage{enumerate}
\usepackage{algorithm}\usepackage{algpseudocode}
\usepackage{tikz}
\newtheorem{corollary}{Corollary}
\newtheorem{lemma}{Lemma}
\DeclareMathOperator*{\sgn}{sgn}
\title{Linear Model Extraction via Factual and Counterfactual Queries}
\author{Daan Otto, Jannis Kurtz, Dick den Hertog, Ilker Birbil}
\date{Source: arXiv:2602.09748v2 (CC BY 4.0)}
\begin{document}
\maketitle

\noindent\emph{Source and licence.} Linear Model Extraction via Factual and Counterfactual Queries. Version arXiv:2602.09748v2 (CC BY 4.0), distributed by arXiv under the Creative Commons Attribution 4.0 International licence, http://creativecommons.org/licenses/by/4.0/, as recorded in the arXiv licence metadata for this exact version and shown on the abstract page https://arxiv.org/abs/2602.09748v2. Full licence notice: \texttt{licenses/arxiv-2602.09748-CC-BY-4.0.txt}.\par\smallskip

\noindent\textit{Editorial scope.} Counted unit: the article's lemma lem:CF\_direction (source0.tex lines 403--409). The article's complete original proof of it is delivered with it (source0.tex lines 843--885, one \texttt{proof} environment of 43 lines, which the article places away from the statement). No mathematics is written, repaired or re-derived: the statement and the proof are the article's own lines, in the article's own notation, and no retained line is substituted or re-flowed.  Retained uncounted as the source-local context the proof needs: lines 378--386.  Nothing was omitted from any retained span, so the package carries no omission record.  No retained line contains a citation, so the wrapper carries no bibliography and no bibliography entry.  No retained line contains a figure environment, an image-inclusion command or drawing code, so no image is delivered and the package declares no asset dependency.  Editorial formatting only, in addition: an AMS \texttt{amsart} preamble that restates the article's own declarations and macros used by the retained lines, namely the environments \texttt{corollary}, \texttt{lemma} on separate counters, together with the standard packages those lines need.  The retained environments print in this document as Corollary 1, Lemma 1 in order of appearance, whereas the article prints the same items with the numbering of its own full document.  The complete native source of the article as distributed in the arXiv e-print for this exact version is bundled unchanged as \texttt{sources/194341/source0.tex}.
\begin{corollary}
\label{cor:optimality_conditions_CF_non_differentiable_norm}
Let $\bm{x}_F\in \mathcal X$ be an arbitrary factual instance and $q_{CF}(\bm{x}_F) = \bm{x}_{CF}^*$ under an arbitrary norm $ \| \cdot \|_{N_1}$.  Then, $\bm{x}_{CF}^*$ is an optimal counterfactual if and only if there exists a $\lambda\in\mathbb R\setminus \{ 0\}$ such that
\begin{align*}
    &\bm{a}^\top \bm{x}_{CF}^* = b, \\
    &  \lambda \bm{a}^\top (\bm{x}_{CF}^* - \bm{x}_F) = \|\bm{x}_{CF}^* - \bm{x}_F\|_{N_1}, \\
    & \|\lambda \bm{a} \|_{N_1}^* \le 1 .
\end{align*}    
\end{corollary}

\begin{lemma}\label{lem:CF_direction}
Let $\bm{x}_F\in \mathcal X$ be an arbitrary factual instance that does not lie on the hyperplane and $q_{CF}(\bm{x}_F) = \bm{x}_{CF}^*$ under an arbitrary norm $ \| \cdot \|_{N_1}$. Then,
\begin{equation}\label{eq:CF_direction}
        \bm{x}_{CF}^*=\bm{x}_F+d_{\bm{x}_F}\bm{v},
    \end{equation}
with $d_{\bm{x}_F}=\frac{b-\bm{a}^{\top}\bm{x}_F}{\|\bm{a}\|_{N_1}^*}$ is an optimal counterfactual if and only if $\|\bm{v}\|_{N_1}\leq 1$ and $\bm{a}^{\top}\bm{v}=\|\bm{a}\|_{N_1}^*$.
\end{lemma}

\begin{proof}[Lemma \ref{lem:CF_direction}]
    We note that $\bm{a}\neq\bm{0}$, hence $\|\bm{a}\|^*\neq\bm{0}$ and $d_{\bm{x}_F}$ is well defined. Besides, since we consider points $\bm{x}_F$ that do not lie on the hyperplane, we know $d_{\bm{x}_F}>0$. Moreover, we have $\|\bm{a}\|_{N_1}^*:=\sup_{\|\bm{x}\|_{N_1}\leq 1}\bm{a}^{\top}\bm{x}$. Since the unit ball defined by $\|\bm{x}\|_{N_1}\leq 1$ is a compact region and the map $\bm{x} \mapsto \bm{a}^{\top}\bm{x}$ is continuous, we know the supremum is attained at a certain point $\bm{v}$ such that $\|\bm{v}\|_{N_1}\leq 1$ and $\bm{a}^{\top}\bm{v}=\|\bm{a}\|_{N_1}^*$.

    First we consider an vector $\bm{v}$ with $\|\bm{v}\|_{N_1}\leq 1$ and $\bm{a}^{\top}\bm{v}=\|\bm{a}\|_{N_1}^*$ and show that $\bm{x}_{CF}^*= \bm{x}_{F}+d_{\bm{x}_F}\bm{v}$ is an optimal counterfactual.
    We use this $\bm{v}$ and $d_{\bm{x}_F}$ to show that the optimality conditions described in Corollary \ref{cor:optimality_conditions_CF_non_differentiable_norm} are met. For the first condition we have
    \begin{align*}
        \bm{a}^\top \bm{x}_{CF}^* &= \bm{a}^\top (\bm{x}_F+d_{\bm{x}_F}\bm{v})= \bm{a}^\top\bm{x}_F +\frac{b-\bm{a}^{\top}\bm{x}_F}{\|\bm{a}\|_{N_1}^*}\bm{a}^\top\bm{v}
        = \bm{a}^\top\bm{x}_F +\frac{b-\bm{a}^{\top}\bm{x}_F}{\|\bm{a}\|_{N_1}^*}\|\bm{a}\|_{N_1}^*\\
        &= \bm{a}^\top\bm{x}_F +b-\bm{a}^{\top}\bm{x}_F=b.
    \end{align*}
    The second and third conditions require a $\lambda\in\mathbb{R}$. Set $\lambda=\frac{\|\bm{v}\|_{N_1}}{\bm{a}^{\top}\bm{v}} \sgn(d_{\bm{x}_F})$. Then, the third condition is met:
    \[ \|\lambda\bm{a}\|_{N_1}^* = |\lambda|   \|\bm{a}\|_{N_1}^* = |\lambda|   \bm{a}^{\top}\bm{v}
    = \frac{\|\bm{v}\|_{N_1}}{|\bm{a}^{\top}\bm{v}|} \bm{a}^{\top}\bm{v}\leq \|\bm{v}\|_{N_1}\leq 1.\]
    Lastly, the second optimality condition is also satisfied:
    \begin{align*}
    \lambda \bm{a}^\top (\bm{x}_{CF}^* - \bm{x}_F) &= \lambda  d_{\bm{x}_F}  \bm{a}^\top \bm{v}
    =\frac{\|\bm{v}\|_{N_1}}{\bm{a}^{\top}\bm{v}}   \sgn(d_{\bm{x}_F})   d_{\bm{x}_F}   \bm{a}^\top \bm{v}\\
    &= |d_{\bm{x}_F}|  \|\bm{v}\|_{N_1} =\|d_{\bm{x}_F}  \bm{v}\|_{N_1} =\|\bm{x}_{CF}^* - \bm{x}_F\|_{N_1} .
    \end{align*}
    We conclude that $\bm{x}_{CF}^*= \bm{x}_{F}+d_{\bm{x}_F}\bm{v}$ is an optimal counterfactual.

    Second, we consider an optimal counterfactual $\bm{x}_{CF}^*= \bm{x}_{F}+d_{\bm{x}_F}\bm{v}$ and will show that $\bm{a}^{\top}\bm{v}=\|\bm{a}\|_{N_1}^*$ with $\|\bm{v}\|_{N_1}\leq 1$. Since $\bm{x}_{CF}^*$ is an optimal counterfactual, Corollary \ref{cor:optimality_conditions_CF_non_differentiable_norm} states that the two optimality conditions are met. The first condition tells us
    \begin{align*}
        \bm{a}^\top \bm{x}_{CF}^* &= \bm{a}^\top (\bm{x}_F+d_{\bm{x}_F}\bm{v})= \bm{a}^\top\bm{x}_F +\frac{b-\bm{a}^{\top}\bm{x}_F}{\|\bm{a}\|_{N_1}^*}\bm{a}^\top\bm{v}=b.
    \end{align*}
    Rearranging the terms yields
    \begin{align*}
        \frac{\bm{a}^\top\bm{v}}{\|\bm{a}\|_{N_1}^*}&=\frac{b-\bm{a}^{\top}\bm{x}_F}{b-\bm{a}^{\top}\bm{x}_F}=1.
    \end{align*}
    Hence, we have $\bm{a}^\top\bm{v}=\|\bm{a}\|_{N_1}^*$. The second optimality condition tells us that there exists a $\lambda\in\mathbb{R}$ such that the following relation holds:
    \begin{align*}
        \lambda \bm{a}^\top (\bm{x}_{CF}^* - \bm{x}_F) =\lambda d_{\bm{x}_F}\bm{a}^\top\bm{v}=\|\bm{x}_{CF}^* - \bm{x}_F\|_{N_1}=\|d_{\bm{x}_F}\bm{v}\|_{N_1}= |d_{\bm{x}_F}|\|\bm{v}\|_{N_1}.
    \end{align*}
    Using the fact that $\bm{a}^\top\bm{v}=\|\bm{a}\|_{N_1}^*$, we can rearrange the terms to find
    \begin{align*}
        \lambda&=\sgn(d_{\bm{x}_F})\frac{\|\bm{v}\|_{N_1}}{\|\bm{a}\|_{N_1}^*}.
    \end{align*}
    The last optimality condition implies a bound on $\|\lambda \bm a\|^*\leq 1$:
    \begin{align*}
        \|\lambda \bm a\|_{N_1}^*& = |\lambda| \|\bm{a}\|_{N_1}^* = |\lambda| \|\bm{a}\|_{N_1}^* = \frac{\|\bm{v}\|_{N_1}}{\|\bm{a}\|_{N_1}^*}\|\bm{a}\|_{N_1}^*=\|\bm{v}\|_{N_1}\leq 1.
    \end{align*}
    Therefore, it must hold that $\|\bm{v}\|_{N_1}\leq 1$ and $\bm{a}^{\top}\bm{v}=\|\bm{a}\|_{N_1}^*$
\end{proof}

\end{document}
